\subsection{Peripheral Extension: Capital Composition and BoP-Constrained Identification}
\label{subsec:theta_identification_chile}

The center identification recovers $\hat{\theta}(\omega_t)$ from a
single cost-minimization problem on a homogeneous capital stock. In
peripheral capitalism, this architecture is insufficient. The
heterogeneity of capital stock imposes an external restriction due to
the import requirements of machinery and equipment ($K^{ME}$), while
assuming that non-residential infrastructure ($K^{NR}$) is domestically
producible but tied to accumulation of machinery and equipment at a
given choice of technique under cost-minimizing principles. Because the
domestic capital-goods sector cannot substitute for imported machinery
at scale, the mechanization decision is constrained by the availability
of foreign exchange. The balance-of-payments constraint enters the
cost-minimization problem through the import content of machinery
investment.

\subsubsection{Capital Decomposition and the Peripheral MPF}
\label{subsec:chile_mpf}

Define the capital stock decomposition:
%
\begin{equation}
	K^{CL} = K^{NR} + K^{ME},
	\qquad
	\phi_t \equiv \frac{K_t^{ME}}{K_t^{CL}} \in (0,1),
	\label{eq:cl:capital_decomp}
\end{equation}
%
where $\phi_t$ is the machinery share of the total non-residential
capital stock. The import content of machinery investment in the early period of interest was registered by 
\citet{Kaldor1959} at $\xi \approx 0.92$--$0.94$. 

Infrastructure and machinery carry differential productivity effects.
Define the composition-weighted MPF slope:
%
\begin{equation}
	\bar{\psi}(\phi)
	=
	\psi^{NR}(1-\phi)+\psi^{ME}\phi,
	\label{eq:cl:psi_bar}
\end{equation}
%
The inequality $\psi^{ME} > \psi^{NR}$ is the structural
content of the \citet{Kaldor1959} hypothesis: embodied technical change is
higher in imported machinery than in domestic infrastructure. Higher
$\phi$ raises productivity potential but deepens import dependence.
The peripheral MPF is:
%
\begin{equation}
	\hat{a}
	=
	\bar{\psi}(\phi)\,\hat{q}
	+
	\lambda_2\,\hat{q}^{\,2},
	\qquad
	\lambda_2 < 0,
	\label{eq:cl:mpf}
\end{equation}
%
Same functional form as the center~\eqref{eq:cost_min} but with a
composition-dependent slope.

\subsubsection{The BoP Penalty and the Effective Profit Share}
\label{subsec:chile_bop_penalty}

Mechanization at rate $\hat{q}$ generates an import bill
$\xi\,\phi\,\hat{q}\,K^{CL}$. Let $\lambda \geq 0$ denote the
shadow cost of foreign exchange --- the scarcity rent under the
BoP constraint. The peripheral cost-minimization problem is:
%
\begin{align}
	\min_{\hat{q}}\; c^{CL}
	&=
	(1-\omega)\hat{q}
	-
	\hat{a}
	+
	\lambda\,\xi\,\phi\,\hat{q}
	\label{eq:cl:cost_min}
	\\[4pt]
	\text{s.t.}\qquad
	\hat{a}
	&=
	\bar{\psi}(\phi)\,\hat{q}
	+
	\lambda_2\,\hat{q}^{\,2},
	\qquad
	\lambda_2<0
	\label{eq:cl:mpf_constraint}
\end{align}

%


The first-order condition yields:
%
\begin{equation}
	\hat{q}^{CL*}(\omega,\phi,\lambda)
	=
	\frac{(1-\omega)-\lambda\,\xi\,\phi-\bar{\psi}(\phi)}
	{2\lambda_2},
	\label{eq:cl:foc_qstar}
\end{equation}
%
and the peripheral transformation elasticity:
%
\begin{equation}
	\boxed{
		\theta^{CL}(\omega,\phi,\lambda)
		=
		\frac{\bar{\psi}(\phi)+\pi^{eff}}{2},
		\qquad
		\pi^{eff}
		\equiv
		(1-\omega)-\lambda\,\xi\,\phi
	}
	\label{eq:cl:theta}
\end{equation}
%
The BoP constraint compresses the effective profit share available for
mechanization. For given $\omega$, the periphery achieves a lower
$\theta$ than the center by exactly $\lambda\xi\phi/2$. The
center solution is recovered as the special case $\lambda = 0$,
$\phi = 0$.

Comparing with the center~\eqref{eq:foc_qstar}:
%
\begin{equation}
	\hat{q}^{US*} - \hat{q}^{CL*}
	=
	\frac{
		\lambda\,\xi\,\phi
		+
		\left[\bar{\psi}(\phi)-\psi_1\right]
	}{
		2|\lambda_2|
	}
	>
	0,
	\label{eq:cl:q_gap}
\end{equation}
%
The gap has two sources: the BoP penalty $\lambda\xi\phi$ and the
structural MPF slope differential
$\bar{\psi}(\phi)-\psi_1$.

\subsubsection{Regime Switch from Complementary Slackness}
\label{subsec:chile_regimes}

The BoP ceiling $\bar{\phi}$ limits the machinery share in new
investment. By complementary slackness, two regimes emerge:
%
\begin{center}
	\begin{tabular}{llll}
		\toprule
		& $\lambda$ & $\phi_{new}$ & $\theta$ \\
		\midrule
		Regime~1 (BoP slack) & $= 0$ & $= \phi^*$ & $\theta^{(1)}$ \\
		Regime~2 (BoP binding) & $> 0$ & $= \bar{\phi} < \phi^*$
		& $\theta^{(2)} < \theta^{(1)}$ \\
		\bottomrule
	\end{tabular}
\end{center}
%
The gap between regimes decomposes as:
%
\begin{equation}
	\theta^{(1)}-\theta^{(2)}
	=
	\underbrace{
		(\psi^{ME}-\psi^{NR})(\phi^*-\bar{\phi})
	}_{\text{composition compression}}
	+
	\underbrace{
		\frac{\lambda\,\xi\,\bar{\phi}}{2}
	}_{\text{forex shadow cost}},
	\label{eq:cl:theta_gap}
\end{equation}
%
The ceiling $\bar{\phi}$ responds to the non-reinvested surplus --- 
profits not converted into machinery investment. \citet{Kaldor1959} 
identifies two competing forces through which this surplus conditions 
the foreign-exchange constraint: lower machinery investment reduces 
current demand for imported capital goods, relaxing the ceiling; yet 
surplus flowing into capitalist consumption tightens it, since luxury 
consumption carries substantially higher import content than average 
consumption once invisible imports are accounted for. Whether the 
accumulation-relief or the consumption-drain mechanism dominates is 
historically contingent and empirically recovered from the long-run 
import propensity relation

\begin{equation}
	m_t
	=
	\zeta_0
	+
	\zeta_1\, k_t^{ME}
	+
	\zeta_2\, nrs_t
	+
	\zeta_3\, \omega_t
	+
	ECT_{m,t},
	\label{eq:cl:imports_coint}
\end{equation}

Where \(m_t\) and  \(nrs_t\) are imports and non-reinvested surplus in log levels, 
such that  $\zeta_1 > 0$ captures the structural import demand generated 
by machinery accumulation and $\zeta_2 \gtrless 0$ recovers the net 
distributional effect without sign restriction, while the $ECT_t$ as repersenting 
structural imbalances would be introuced below as a regime switching variable for 
the identification of sudden shifts in the dynamics of unbalanced growth. 

\subsubsection{First Stage: Cointegration System for Propensity to Import}
\label{subsec:chile_imports}

The first stage identifies the long-run external relation through a
cointegration system to measure import $M$, propensities from $NRS$ non-reinvested surplus, $\omega$ wage share, and ${K_t}^ME$ represented by the following state vector in log levels and the wage share:

\begin{equation}
	Y_t
	=
	\left(
	m_t,\;
	k_t^{ME},\;
	nrs_t,\;
	\omega_t,
	\right)',
	\label{eq:cl:state_vector_theta}
\end{equation}


Theoretically, $\zeta_1 > 0$ captures the extent to which
non-reinvested surplus pressures the foreign-exchange constraint, while
$\zeta_2 > 0$ captures the structural import demand generated by
machinery accumulation. The estimated parameters and their relative magnitude, 
would be suggestive of the effects of distribution $\zeta_3$ \citep{Kaldor1959}. 
The estimated disequilibrium terms of the co-integrated system, would 
summarize the externally driven systemic imbalances driven the scarcity 
of foreign exchange and their use to import machinery and equipment 
relative to imported luxury consumption driven by the non-reinvested surplus.  

\subsubsection{Second Stage: Threshold VECM for the Productive Frontier}
\label{subsec:chile_frontier_threshold}

The state vector for the productive frontier is:
%
\begin{equation}
	X_t^{CL,\theta}
	=
	\left(
	y_t,\;
	k_t^{NR},\;
	k_t^{ME},\;
	\omega_t k_t^{ME}
	\right)',
	\label{eq:cl:state_vector_theta}
\end{equation}
%
The interaction term is $\omega_t k_t^{ME}$ because the
cost-minimization problem shows that distribution conditions technique
choice through the machinery component specifically. Infrastructure is
distribution-insensitive at business-cycle frequency.

The first cointegrating vector identifies the capacity frontier:
%
\begin{equation}
	y_t
	=
	\kappa_1
	+
	\theta_0\, k_t^{NR}
	+
	\psi\, k_t^{ME}
	+
	\theta_2\,(\omega_t k_t^{ME})
	+
	ECT_{\theta,t},
	\label{eq:cl:cv1}
\end{equation}
%
where $\psi = \psi^{ME}-\psi^{NR}$ is the machinery
productivity premium and $\theta_2 < 0$ is the
distribution-composition interaction.

This frontier relation does not directly identify the level of capacity
utilization. The term $ECT_{\theta,t}$ is a long-run disequilibrium term
of the frontier system. It identifies deviations between realized
output and the productive-capacity frontier implied by capital
composition and distributional conditions, but it does not, by itself,
identify the level of $\mu_t^{CL}$.

The empirical transformation elasticity is instead recovered from the
frontier coefficients as a composition-sensitive object:
%
\begin{equation}
	\hat{\theta}^{CL}(\omega_t,\phi_t)
	=
	\theta_0
	+
	\psi\,\phi_t
	+
	\theta_2\,\omega_t\,\phi_t,
	\label{eq:cl:theta_hat}
\end{equation}
%
which is the empirical counterpart of
\eqref{eq:cl:theta}.

To identify nonlinear external activation, the productive frontier is
estimated as a threshold VECM in the sense of \citet{Seo2006}, using
the first-stage external disequilibrium as the transition object. Let
the regime indicator be
%
\begin{equation}
	R_t
	=
	\mathbf{1}\!\left[\widehat{ECT}_{m,t-1} > \lambda\right],
	\label{eq:cl:regime_indicator}
\end{equation}
%
where $\lambda$ is interpreted as the threshold that activates the
external constraint as a shadow price.

The threshold VECM is then
%
\begin{equation}
	\Delta X_t^{CL,\theta}
	=
	\mu
	+
	\alpha^{(1)} ECT_{\theta,t-1}(1-R_t)
	+
	\alpha^{(2)} ECT_{\theta,t-1}R_t
	+
	\sum_{j=1}^{p}\Gamma_j \Delta X_{t-j}^{CL,\theta}
	+
	\varepsilon_t,
	\label{eq:cl:tvecm_theta}
\end{equation}
%
where the threshold effect is driven by the estimated external
disequilibrium recovered from the first-stage import system. In this
architecture, the threshold is not imposed on the import-propensity
equation itself. Rather, the import system identifies the external
imbalance, and the threshold VECM identifies how that imbalance
activates the shadow-price effect on the machinery-sensitive productive
frontier.

Productive-capacity growth is then constructed through the closure
%
\begin{equation}
	g_{Y_t^{p,CL}}
	=
	\hat{\theta}^{CL}(\omega_t,\phi_t)\,g_{K_t^{CL}},
	\label{eq:cl:frontier_growth}
\end{equation}
%
so that the method identifies the evolution of utilization only in
growth rates:
%
\begin{equation}
	g_{\hat{\mu}_t^{CL}}
	=
	g_{Y_t^{CL}}
	-
	g_{Y_t^{p,CL}},
	\label{eq:cl:mu_growth}
\end{equation}
%
The level of $\hat{\mu}_t^{CL}$ is recovered only after imposing a
pin-year normalization:
%
\begin{equation}
	\hat{\mu}_{t_0}^{CL}
	=
	\mu_0,
	\label{eq:cl:mu_pin}
\end{equation}
%
and accumulating the productive-capacity series forward and backward
from $t_0$.

\subsubsection{Recursive Identification}
\label{subsec:chile_recursive}

The recursive architecture is therefore clear. In a first stage, a
cointegration system for propensity to import is estimated, and its
estimated parameters and estimated disequilibrium term summarize the
externally driven systemic imbalances. In a second stage, a threshold
VECM for the productive frontier is estimated, using those first-stage
estimations to interpret $\lambda$ as the threshold that activates the
external constraint as a shadow price on the machinery-sensitive
transformation of accumulation into productive capacity. The frontier
system therefore identifies the evolution of capacity utilization only
up to a normalization, while the threshold activation identifies the
foreign-exchange penalty that compresses the feasible mechanization
path. The two systems are brought together in interpretation through the
theoretical object $\theta^{CL}(\omega_t,\phi_t,\lambda)$.

